% =============================================================================
%  Candidate chemical-formula generation for three material systems
%  Methodology, datasets, filtering and diversity report
%
%  Compile with:  pdflatex CANDIDATE_GENERATION_REPORT.tex   (twice for refs)
% =============================================================================
\documentclass[11pt,a4paper]{article}

\usepackage[T1]{fontenc}
\usepackage{lmodern}
\usepackage[margin=2.6cm]{geometry}
\usepackage{amsmath,amssymb}
\usepackage{booktabs}
\usepackage{array}
\usepackage{enumitem}
\usepackage{caption}
\usepackage{xcolor}
\usepackage[detect-all]{siunitx}
\usepackage[colorlinks=true,linkcolor=blue!60!black,citecolor=blue!60!black,
            urlcolor=blue!60!black]{hyperref}

\captionsetup[table]{skip=4pt,font=small,labelfont=bf}
\setlist{nosep}
\setlength{\parskip}{4pt}
\setlength{\parindent}{0pt}
\setlength{\tabcolsep}{4pt}
\emergencystretch=1.5em

\newcommand{\pkg}[1]{\texttt{#1}}
\newcommand{\code}[1]{\texttt{#1}}
\newcommand{\cmd}[1]{\texttt{\small #1}}

\title{\textbf{Candidate Chemical-Formula Generation for Three Material Systems}\\[4pt]
       \large Methodology, reference datasets, novelty filtering and diversity selection}
\author{Chemical Formula Parser\\
        \small\texttt{chemical-formula-parser} (\code{Chemical-Formula-Parser})}
\date{21 September 2026}

\begin{document}
\maketitle

\begin{abstract}
\noindent
We describe the pipeline used to generate 1000 candidate chemical formulas for each of three
composition spaces --- oxychalcogenides, phosphorus chalcogenides and double perovskites.
Candidates are produced by an exact, deterministic enumeration of charge-balanced integer
stoichiometries over a \code{config.yaml}-declared element universe, screened for chemical
plausibility with SMACT~4.0.0 (charge neutrality under the ICSD-2024 oxidation-state set plus
the Pauling electronegativity test), filtered against the union of the Crystallography Open
Database, the Materials Project and the Open Quantum Materials Database
(3{,}597{,}064 unique known reduced formulas), and reduced to a maximally spread-out subset by
greedy farthest-point (max--min dispersion) sampling over a combined stoichiometric and
element-identity distance. All reported numbers are measured from the delivered runs. The three
novelty-filtered outputs contain 1000 formulas each, with zero duplicates, zero parenthesised
groups, zero SMACT-invalid compositions and zero overlap with the reference databases.
\end{abstract}

\tableofcontents
\vspace{2em}

% =============================================================================
\section{Summary}
% =============================================================================

For each composition space the pipeline has four stages:

\begin{enumerate}[label=(\arabic*)]
  \item \textbf{Enumerate} charge-balanced compositions from the config-declared element universe
        (cation set $\times$ anion set $\times$ integer stoichiometry), stratified so that no
        single element combination dominates the pool.
  \item \textbf{Screen} every candidate with SMACT~\cite{smact_joss,smact_chem}
        (\code{smact.screening.smact\_validity}): charge neutrality under an oxidation-state set
        \emph{and} the Pauling electronegativity rule~\cite{pauling}.
  \item \textbf{Filter} out any formula whose canonical reduced form already occurs in
        COD~\cite{cod}, MP~\cite{mp_apl,mp_api} or OQMD~\cite{oqmd}.
  \item \textbf{Select} a maximally spread-out subset of the survivors with greedy farthest-point
        sampling~\cite{gonzalez} over a combined L2 / Jaccard~\cite{jaccard} distance.
\end{enumerate}

Output is written one formula per line in flat, parenthesis-free notation
(\code{La4S3O3}, not \code{La4(SO)3}).

\begin{table}[h]
\centering
\footnotesize
\begin{tabular}{@{}lrrrrrr@{}}
\toprule
System & Pool & SMACT rej. & Known in DBs & Novel & Elem.\ sets & Min $d$\\
\midrule
Oxychalcogenides         & 14{,}975 &  0 (0.0\,\%) &   830 (5.5\,\%) & 14{,}145 & 1{,}000 & 0.5354\\
Phosphorus chalcogenides &  6{,}768 & 36 (0.5\,\%) &    81 (1.2\,\%) &  6{,}651 &   898   & 0.2720\\
Double perovskites       & 11{,}019 &  0 (0.0\,\%) & 2{,}088 (18.9\,\%) & 8{,}931 & 1{,}000 & 0.4905\\
\bottomrule
\end{tabular}
\caption{Statistics of the three delivered runs ($n=1000$, novelty filter on, diversity selection on).
``Min $d$'' is the minimum pairwise distance of the delivered selection.}
\label{tab:runs}
\end{table}

Verification performed on the six output files: 1000 lines each; 0 parenthesised groups;
0 duplicate formulas; 0 formulas failing the SMACT validity test; and, for the three
\code{novel\_1000.txt} files, 0 canonical forms present in the 3{,}597{,}064 known formulas.
The three \code{candidates\_1000.txt} comparison sets are diversity-selected but deliberately
\emph{not} novelty-filtered, so they do contain known formulas (70, 16 and 115 respectively).

% =============================================================================
\section{The three composition spaces}
% =============================================================================

The chemistry of each system is declared entirely in \code{config.yaml} under
\code{composition\_spaces}; no element list is hard-coded in Python. Every space specifies a
cation pool, an anion pool, per-element oxidation states used for charge balancing, and
enumeration constraints (Table~\ref{tab:spaces}).

\begin{table}[h]
\centering
\footnotesize
\begin{tabular}{@{}p{3.5cm} p{4.9cm} p{2.0cm} p{3.4cm}@{}}
\toprule
Space & Cations & Anions & Constraints\\
\midrule
\code{oxychalcogenides} &
La, Ce, Pr, Nd, Sm, Eu, Gd, Ti, V, Cr, Mn, Fe, Co, Ni (14) &
O, S, Se, Te; O required &
max\_stoich 8, max\_total\_atoms 20, max\_cation\_species 3, per\_combo\_limit 4\\
\addlinespace
\code{phosphorus\_\allowbreak chalcogenides} &
Ti, V, Cr, Mn, Fe, Co, Ni, Cu (8) &
P, S, Se, Te; P required &
max\_stoich 8, max\_total\_atoms 20, max\_cation\_species 4, per\_combo\_limit 6\\
\addlinespace
\code{double\_perovskites} &
La, Nd, Sm, Eu, Gd, Ca, Sr, Ba, Ti, V, Cr, Mn, Fe, Co, Ni, Cu, Nb, Mo, W (19) &
O &
max\_stoich 6, max\_total\_atoms 10, max\_cation\_species 4, per\_combo\_limit 6\\
\bottomrule
\end{tabular}
\caption{Composition spaces as declared in \code{config.yaml}.}
\label{tab:spaces}
\end{table}

Chemical rationale for the caps: rare-earth/transition-metal oxychalcogenides (photocatalysis,
thermoelectrics, ion storage) are mostly reported with at most three distinct cation species;
layered transition-metal phosphorus chalcogenides (e.g.\ \code{TiPS3}, \code{FeCoPS}) are
one-to-four cation systems; and the $A_2BB'\mathrm{O}_6$ double-perovskite motif contains at most
four distinct cation species and at most ten atoms per formula unit.

\code{max\_cation\_species} matters for tractability as well as chemistry: without it the
double-perovskite space walks all $2^{19}$ cation subsets and one run took \SI{272}{s}; capped at
four the same run takes \SI{5}{s}.

% =============================================================================
\section{Candidate generation algorithm}
% =============================================================================

Implementation: \pkg{src/composition.py} (\code{CompositionSpace},
\code{generate\_charge\_balanced\_\allowbreak stoichiometries},
\code{check\_charge\_balance}).

For each pair (cation subset $C$, anion subset $A$) satisfying
\[
\begin{aligned}
  &\text{min\_cations} \le |C| \le \text{max\_cation\_species},\\
  &\text{min\_anions} \le |A| \le \text{max\_anion\_species},\\
  &\text{required\_anions} \subseteq A,
\end{aligned}
\]
the pipeline performs the following steps:

\begin{enumerate}[label=\arabic*.]
  \item Enumerate integer stoichiometry vectors over $C \cup A$ with each coefficient between
        $1$ and \code{max\_stoich} and total atom count at most \code{max\_total\_atoms}. A work
        budget (\code{max\_scan}, default 20{,}000 vectors per combination) bounds the cost of
        large combinations.
  \item For each vector, test whether an assignment of oxidation states --- taken from the
        per-element lists in \code{config.yaml} --- makes the composition exactly charge neutral.
        Only vectors admitting such an assignment survive.
  \item Canonicalise the surviving composition with \code{formula\_utils.normalize\_formula}
        (pymatgen~\cite{pymatgen} \code{Composition.get\_reduced\_formula} semantics) and
        deduplicate, so that \code{AlCrNb2} and \code{Nb2CrAl} collapse to one entry.
  \item Keep at most \code{per\_combo\_limit} formulas per $(C,A)$ combination (stratification,
        Section~\ref{sec:strat}) and stop at \code{max\_results} overall.
\end{enumerate}

The enumeration is exact and deterministic: it neither samples randomly nor depends on dictionary
ordering, so re-running a script reproduces byte-identical output.

% =============================================================================
\section{Chemical-plausibility screening with SMACT}
% =============================================================================

Library: SMACT (\emph{Semiconducting Materials by Analogy and Chemical Theory}), version 4.0.0
\cite{smact_joss,smact_chem}. Wrapper: \pkg{src/smactlib/filters.py}; CLI wiring:
\pkg{src/candidate\_cli.py} (\code{\_smact\_validator}, \code{\_screen}).

Every enumerated candidate is passed through
\begin{verbatim}
    smact.screening.smact_validity(formula,
                                   use_pauling_test=True,
                                   include_alloys=True,
                                   oxidation_states_set="icsd24")
\end{verbatim}
so a candidate is kept only if SMACT can find an oxidation-state assignment, drawn from the
ICSD-derived 2024 oxidation-state table shipped with SMACT, that
\begin{enumerate}[label=(\alph*)]
  \item sums to zero charge for the given stoichiometry, and
  \item obeys Pauling's electronegativity rule~\cite{pauling}: negative oxidation states are
        assigned only to the more electronegative elements of the composition.
\end{enumerate}
All-metallic compositions are auto-accepted as alloys (\code{include\_alloys=True}), which is
SMACT's convention for intermetallics where integer charge balancing is not meaningful.

Because the \code{config.yaml} enumeration already enforces charge neutrality with its own
oxidation-state lists, SMACT acts here mainly as an independent cross-check that additionally
imposes the electronegativity constraint. The measured rejection rates confirm that the two are
largely consistent:
\[
  \text{oxychalcogenides } \frac{0}{14{,}975}, \qquad
  \text{P-chalcogenides } \frac{36}{6{,}768}\ (0.5\%), \qquad
  \text{double perovskites } \frac{0}{11{,}019}.
\]

The 36 rejected phosphorus-chalcogenide candidates are Te-rich compositions such as
\code{TiTe4P}, \code{VTe5P}, \code{CrTe7P2} and \code{Cr(Te4P)2}. Pauling electronegativities put
P (2.19) marginally above Te (2.10), so the only charge-balancing assignment for these formulas
requires Te to carry the negative state and P the positive one, which the Pauling test forbids.
Screening costs \SI{0.8}{ms} per formula and is applied to the whole pool, i.e.\ a few seconds
per run. It is on by default and can be retuned from the command line: \cmd{--no-smact},
\cmd{--smact-ox-set \{icsd24,icsd16,smact14,pymatgen\_sp,wiki\}}, \cmd{--smact-no-en}.

\subsection*{SMACT as an independent generator (complementary path)}

\pkg{src/smactlib/generator.py} also implements full SMACT-driven candidate generation: for every
element combination of the requested size it calls \code{smact.neutral\_ratios} over the SMACT
oxidation-state sets, optionally applying \code{smact.screening.pauling\_test}, and enumerates
integer stoichiometry vectors for all-metal (alloy) combinations. That path is driven by
\cmd{scripts/generate\_formulas.py} (which writes \code{data/formulas/candidates.csv} with
per-candidate SMACT flags) and screened for novelty by \cmd{scripts/find\_novel\_formulas.py}.
It searches an arbitrary element universe rather than a cation/anion-structured space, and is
therefore complementary to the three system-specific pipelines. The 3{,}000 delivered formulas
were produced by the config-driven enumerator of Section~3 and validated by the SMACT screen
described above.

% =============================================================================
\section{Reference datasets used for novelty filtering}
% =============================================================================

Three databases were used, each mirrored locally under \code{\textasciitilde/Git/<name>-Dataset}
and reduced to a list of unique canonical formulas in \code{datasets/} by
\cmd{scripts/extract\_database\_formulas.py}. That extraction is run once; the CSVs are the
reusable artefact, and \code{datasets/} is intentionally \code{.gitignore}d because of its size.

\begin{table}[h]
\centering
\footnotesize
\begin{tabular}{@{}l p{7.4cm} r r@{}}
\toprule
Database & Source snapshot & Unique formulas & Size\\
\midrule
COD~\cite{cod} & rsync mirror of \code{crystallography.net/cif} (experimental crystal structures) & 399{,}556 & \SI{6.4}{MB}\\
MP~\cite{mp_apl,mp_api} & \code{data/mp\_summary.csv} via the Materials Project API (DFT-computed entries) & 104{,}734 & \SI{1.1}{MB}\\
OQMD~\cite{oqmd} & SQL dump \code{qmdb\_v1\_8}, release 2026-01-29, \code{compositions} table (DFT-computed entries) & 3{,}215{,}729 & \SI{35.4}{MB}\\
\midrule
 & Sum & 3{,}720{,}019 & \\
 & Union (deduplicated) & 3{,}597{,}064 & \\
\bottomrule
\end{tabular}
\caption{Reference databases and the extracted known-formula sets.}
\label{tab:datasets}
\end{table}

Cross-database overlap within the union: COD $\cap$ MP $= 25{,}971$; COD $\cap$ OQMD $= 47{,}141$;
MP $\cap$ OQMD $= 75{,}399$; present in all three $= 25{,}556$. Hence 122{,}955 formulas are
shared by at least two databases and the union is $3.3\,\%$ smaller than the sum.

Composition of the union by number of distinct elements: 1 element, 104; 2 elements, 185{,}555;
3 elements, 842{,}541; 4 elements, 1{,}973{,}133; 5 elements, 395{,}768; 6 elements, 128{,}402;
7 or more, $\approx$71{,}000.

Extraction details worth recording:
\begin{itemize}
  \item Every formula is re-canonicalised with \code{normalize\_formula} before being written, so
        the three CSVs share one notation and plain string comparison is a valid novelty test.
        Canonicalisation was verified to be notation- and element-order-independent (200 random
        element-order permutations per sample formula, 0 inconsistencies).
  \item COD formulas come from parsing the CIF mirror, MP from the summary CSV, and OQMD from the
        SQL dump. The OQMD dump packs many rows into a single \code{INSERT} statement, so
        extraction matches every \code{('formula',} tuple on \code{compositions} lines rather than
        only the first one per line; an early version that matched only the first recovered 181
        formulas instead of 3.2\,M.
  \item A small number of OQMD rows use placeholder site labels instead of element symbols. These
        fail element parsing and are dropped. In the unlikely case that a placeholder-only row used
        real element symbols it would survive as a spurious ``known'' entry; that can only cause
        false negatives (a genuinely novel formula reported as known), never false positives, and
        no such case was observed for the three heavy-cation systems studied here.
\end{itemize}

\noindent Regenerate the CSVs with:
\begin{verbatim}
    conda run -n chem python scripts/extract_database_formulas.py
    conda run -n chem python scripts/extract_database_formulas.py --db cod mp
\end{verbatim}

% =============================================================================
\section{How novelty filtering is done}
% =============================================================================

Implementation: \pkg{src/novelty\_filter.py} (\code{load\_known\_\allowbreak formulas\_set},
\code{generate\_novel\_\allowbreak candidates}, \code{NovelGenerationResult}).

\begin{enumerate}[label=\arabic*.]
  \item Load the union of the selected CSVs into a Python \code{set} of canonical strings
        (\cmd{--db cod mp oqmd} selects a subset; the default is all three). Loading 3{,}597{,}064
        formulas takes a few seconds and dominates runtime.
  \item Canonicalise each generated candidate with \code{normalize\_formula} and test set
        membership. A candidate is \emph{novel} if and only if its exact reduced formula is absent
        from every selected database.
  \item In enumeration-order mode the generation limit is scaled adaptively ($\times 4$ per round,
        capped at 2{,}000{,}000) until the requested number of novel formulas is found, the
        composition space is exhausted, or the cap is hit. The result object reports
        \code{total\_generated}, \code{rejected\_by\_validator}, \code{known\_filtered},
        \code{novel\_found}, \code{space\_exhausted} and \code{target\_met}, so statistics are
        never inferred from a truncated list.
\end{enumerate}

Novelty rates observed at pool level (fraction of charge-balanced, SMACT-valid candidates that no
database contains) are given in Table~\ref{tab:runs}: $94.5\,\%$ for oxychalcogenides, $98.8\,\%$
for phosphorus chalcogenides and $81.1\,\%$ for double perovskites. The lower double-perovskite
rate is expected --- oxide perovskite compositions are the most heavily mined region of all three
databases.

\textbf{Scope.} ``Novel'' here means \emph{database-unseen at the level of the reduced chemical
formula}. It is not a claim of thermodynamic stability, synthesizability, or that the material has
never been made. The comparison is also composition-level only: two polymorphs of the same
composition count as one known formula.

% =============================================================================
\section{How diversity is achieved}
% =============================================================================

Diversity is produced by two cooperating mechanisms: stratified enumeration, which makes the pool
broad, and farthest-point selection, which makes the delivered subset spread out.

\subsection{Stratified enumeration (pool construction)}
\label{sec:strat}

Raw enumeration exhausts every stoichiometry of one element combination before moving to the next,
and iterates cation subsets by increasing size, so the head of an unstratified list is dominated by
a handful of species: in the first oxychalcogenide runs the first 1000 formulas spanned only 15
distinct element sets and were nearly all single-cation La systems. \code{per\_combo\_limit} caps
how many formulas any single (cation-set, anion-set) combination contributes (4 for
oxychalcogenides, 6 for the other two systems), spreading the pool across the whole combination
space.

\subsection{Distance metric}

Implementation: \pkg{src/diversity.py}. Each formula is mapped to two 118-dimensional vectors over
the periodic table: $\mathbf{f}$, the elemental mole fractions (summing to 1), and $\mathbf{p}$,
the binary element-presence indicator. The distance between two formulas is
\begin{equation}
  d(i,j) \;=\; w_{\text{comp}}\,\bigl\lVert \mathbf{f}_i - \mathbf{f}_j \bigr\rVert_2
           \;+\; w_{\text{set}}\,\Bigl(1 - J(S_i,S_j)\Bigr),
  \qquad
  J(S_i,S_j) = \frac{|S_i \cap S_j|}{|S_i \cup S_j|},
  \label{eq:distance}
\end{equation}
with $J$ the Jaccard index~\cite{jaccard} on element sets. The first term separates
stoichiometries; the second separates chemical identity, so two formulas sharing no elements are
pushed apart even when their fraction vectors look similar. Both weights default to 1.0 and are
exposed as \cmd{--w-composition} and \cmd{--w-element-set} (e.g.\ \cmd{--w-element-set 3}
emphasises distinct species over distinct stoichiometry).

\subsection{Selection algorithm: greedy farthest-point sampling}

\code{select\_diverse(pool, n\_select)} implements max--min dispersion~\cite{gonzalez}:
\begin{itemize}
  \item deduplicate the pool and seed with the candidate farthest from the pool centroid (a
        deterministic, representative starting point);
  \item maintain $m_i = \min_{s \in \text{selected}} d(i,s)$ for every pool member $i$;
  \item repeatedly select $\arg\max_i m_i$ and update
        $\mathbf{m} \leftarrow \min\bigl(\mathbf{m},\, d(\cdot, i^\ast)\bigr)$, marking selected
        entries as $-\infty$.
\end{itemize}
This is Gonzalez's 2-approximation for the max--min dispersion ($p$-center) objective~\cite{gonzalez}:
the smallest pairwise distance in the returned set is at least half of the best achievable for that
set size. Vectorised with NumPy it costs $O(|\text{pool}| \times n_{\text{select}})$ distance
evaluations --- reducing a 15{,}000-member pool to 1000 formulas takes a few seconds. The result is
deterministic and reproducible.

\subsection{Measured effect}

\begin{table}[h]
\centering
\small
\begin{tabular}{@{}lccc@{}}
\toprule
 & Baseline A & Baseline B & Delivered\\
 & no stratification, & stratified, & stratified $+$ FPS\\
 & enumeration order & enumeration order & $+$ novelty filter\\
\midrule
Oxychalcogenides: distinct sets &   15 &  258 & 1{,}000\\
\quad min pairwise distance     & 0.0157 & 0.0288 & 0.5354\\
\quad mean pairwise distance    & ---    & 1.2574 & 1.2975\\
\addlinespace
P-chalcogenides: distinct sets  &    7 &  167 &   898\\
\quad min pairwise distance     & 0.0157 & 0.0315 & 0.2720\\
\quad mean pairwise distance    & ---    & 1.1824 & 1.1616\\
\addlinespace
Double perovskites: distinct sets &  155 &  228 & 1{,}000\\
\quad min pairwise distance     & 0.0354 & 0.0354 & 0.4905\\
\quad mean pairwise distance    & ---    & 1.1544 & 1.0689\\
\bottomrule
\end{tabular}
\caption{Diversity of 1000-formula outputs. Baseline~A was measured on the pre-stratification
pipeline; Baseline~B is reproducible today with \cmd{--no-diverse}; ``Delivered'' is the content of
the \code{novel\_1000.txt} files.}
\label{tab:diversity}
\end{table}

Element coverage is complete in all delivered sets (18, 12 and 20 elements respectively --- every
element of the space appears), and the largest group of formulas sharing one element set drops from
4--6 to 1--3. Phosphorus chalcogenides reach 898 rather than 1000 distinct element sets because the
space contains only 12 elements with P mandatory, so fewer than 1000 distinct sets are reachable at
this size.

\textbf{Honest caveat.} Max--min dispersion optimises the \emph{minimum} pairwise distance, not the
mean. The mean can therefore fall slightly when the selection is forced to fill in isolated
regions: for double perovskites it went from 1.1544 (Baseline~B) to 1.0689 (delivered) while the
minimum rose from 0.0354 to 0.4905 and distinct element sets rose from 228 to 1000. If average
separation is the target instead, \cmd{--no-diverse} with a larger pool, or a max--sum objective,
would be the right tool.

% =============================================================================
\section{Output notation}
% =============================================================================

\code{formula\_utils.flatten\_formula} expands parenthesised groups and re-emits the reduced
composition with integer counts, keeping pymatgen's~\cite{pymatgen} conventional
electronegativity-based element order (cation before anion): \code{La4(SO)3} becomes
\code{La4S3O3}, \code{Ti(PS3)2} becomes \code{TiP2S6}. Unit counts are omitted
(\code{LaTiSO2}, not \code{La1Ti1S1O2}). Flattening is notation-independent, idempotent and
round-trips to the same canonical form as \code{normalize\_formula} (covered by tests in
\pkg{tests/test\_formula\_utils.py}), so flat output can still be matched against the parenthesised
database CSVs by canonicalising both sides --- which is exactly how the zero-overlap verification
in Section~1 was performed.

% =============================================================================
\section{Reproducibility}
% =============================================================================

Each system has its own script; all share \pkg{src/candidate\_cli.py}.

{\footnotesize
\begin{verbatim}
    conda run -n chem python examples/oxychalcogenides/generate_candidates.py \
        -n 1000 --filter-known -o examples/oxychalcogenides/novel_1000.txt

    conda run -n chem python examples/phosphorus_chalcogenides/generate_candidates.py \
        -n 1000 --filter-known -o examples/phosphorus_chalcogenides/novel_1000.txt

    conda run -n chem python examples/double_perovskites/generate_candidates.py \
        -n 1000 --filter-known -o examples/double_perovskites/novel_1000.txt
\end{verbatim}
\par}

{\sloppy
The comparison sets drop \cmd{--filter-known} and write
\code{candidates\_\allowbreak 1000.txt}.
\cmd{examples/\allowbreak generate\_candidates.py} runs all three spaces in one go, inserting the
space name into the output filename.

\textbf{Flags.} \cmd{-n/--num}, \cmd{-o/--output}, \cmd{--filter-known}, \cmd{--db cod mp oqmd},
\cmd{--no-diverse}, \cmd{--pool}, \cmd{--per-combo}, \cmd{--max-cations}, \cmd{--max-anions},
\cmd{--w-composition}, \cmd{--w-element-set}, \cmd{--no-smact}, \cmd{--smact-ox-set},
\cmd{--smact-no-en}.
\par}

\textbf{Cost.} A full 1000-formula run (pool construction $+$ SMACT screen $+$ loading the
3.6\,M-formula known index $+$ FPS selection) takes \SI{17}{s} wall clock; without
\cmd{--filter-known} it is a few seconds. Output is deterministic.

\textbf{Tests.} \cmd{conda run -n chem python -m pytest -q} collects 100 tests, all passing,
covering formula canonicalisation and flattening, charge-balance enumeration, the novelty and
validator hooks, SMACT filtering and the diversity selector.

% =============================================================================
\section{Limitations and scientific safeguards}
% =============================================================================

\begin{itemize}
  \item Charge neutrality plus the Pauling electronegativity rule are \emph{necessary}, not
        sufficient, conditions for a real compound. They screen out impossible stoichiometries; they
        say nothing about thermodynamic stability, synthesis conditions or whether a phase can be
        isolated. No DFT or machine-learning stability prediction (e.g.\ formation energy above the
        convex hull) is applied at this stage.
  \item ``Novel'' means unseen in COD, MP and OQMD at the reduced-formula level. Absence from these
        databases is evidence of interest, not evidence of non-existence: the databases are
        incomplete, and unpublished or very recent compounds will be counted as novel.
  \item Oxidation states are taken from \code{config.yaml} for enumeration and from SMACT's
        \code{icsd24} set for screening. Systems with unusual or mixed-valence chemistry may need a
        different set (\cmd{--smact-ox-set}) or a config edit.
  \item Diversity is measured in composition space only. Structurally distinct polymorphs of one
        composition are a single point; conversely, two compositions far apart in metric
        \eqref{eq:distance} may still be structurally related.
  \item Per-combination work budgets (\code{max\_scan}) and species caps
        (\code{max\_cation\_species}) mean the pool is a large, stratified sample of the space
        rather than a proof of exhaustive coverage.
\end{itemize}

% =============================================================================
\section*{Software availability}
% =============================================================================

Pipeline code: \code{chemical-formula-parser} at
\code{/home/huangyang/Git/Chemical-Formula-Parser} (modules \pkg{composition}, \pkg{diversity},
\pkg{novelty\_filter}, \pkg{formula\_utils}, \pkg{candidate\_cli}, \pkg{smactlib}).
Environment: conda env \code{chem} with SMACT~4.0.0, pymatgen, NumPy, SciPy, pandas, PyYAML.
Reference-database extracts: \code{datasets/\{cod,mp,oqmd\}\_formulas.csv}, regenerated by
\cmd{scripts/extract\_database\_formulas.py}.

% =============================================================================
\begin{thebibliography}{10}
\providecommand{\doi}[1]{doi:~\href{https://doi.org/#1}{\nolinkurl{#1}}}

\bibitem{cod}
S.~Gra\v{z}ulis, D.~Chateigner, R.~T.~Downs, A.~F.~T.~Yokochi, M.~Qui\'{r}os, L.~Lutterotti,
E.~Manakova, J.~Butkus, P.~Moeck, and A.~Le~Bail.
\newblock Crystallography Open Database --- an open-access collection of crystal structures.
\newblock \emph{Journal of Applied Crystallography}, 42(4):726--729, 2009.
\newblock \doi{10.1107/S0021889809016690}

\bibitem{mp_apl}
A.~Jain, S.~P.~Ong, G.~Hautier, W.~Chen, W.~D.~Richards, S.~Dacek, S.~Cholia, D.~Gunter,
D.~Skinner, G.~Ceder, and K.~A.~Persson.
\newblock Commentary: The Materials Project: A materials genome approach to accelerating materials
innovation.
\newblock \emph{APL Materials}, 1(1):011002, 2013.
\newblock \doi{10.1063/1.4812323}

\bibitem{mp_api}
Materials Project.
\newblock Data accessed via the MP API; local snapshot \code{mp\_summary.csv}
(\code{\textasciitilde/Git/MP-Dataset}).

\bibitem{oqmd}
S.~Kirklin, J.~E.~Saal, B.~Meredig, A.~Thompson, J.~W.~Doak, M.~Aykol, S.~R\"{u}hl, and
C.~Wolverton.
\newblock The Open Quantum Materials Database (OQMD): assessing the accuracy of DFT formation
energies.
\newblock \emph{npj Computational Materials}, 1:15010, 2015.
\newblock \doi{10.1038/npjcompumats.2015.10}
\newblock Snapshot used: QMDB v1.8 SQL dump, release 2026-01-29.

\bibitem{smact_joss}
D.~W.~Davies, K.~T.~Butler, A.~J.~Jackson, J.~M.~Skelton, K.~Morita, and A.~Walsh.
\newblock SMACT: Semiconducting Materials by Analogy and Chemical Theory.
\newblock \emph{Journal of Open Source Software}, 4(38):1361, 2019.
\newblock \doi{10.21105/joss.01361}
\newblock Version used: 4.0.0.

\bibitem{smact_chem}
D.~W.~Davies, K.~T.~Butler, A.~J.~Jackson, A.~Morris, J.~M.~Frost, J.~M.~Skelton, and A.~Walsh.
\newblock Computational Screening of All Stoichiometric Inorganic Materials.
\newblock \emph{Chem}, 1(4):617--627, 2016.
\newblock \doi{10.1016/j.chempr.2016.09.010}

\bibitem{pauling}
L.~Pauling.
\newblock \emph{The Nature of the Chemical Bond}.
\newblock Cornell University Press, Ithaca, NY, 3rd edition, 1960.
\newblock Electronegativity criterion implemented by \code{smact.screening.pauling\_test}.

\bibitem{gonzalez}
T.~F.~Gonzalez.
\newblock Clustering to minimize the maximum intercluster distance.
\newblock \emph{Theoretical Computer Science}, 38:293--306, 1985.
\newblock \doi{10.1016/0304-3975(85)90224-5}
\newblock Greedy farthest-point sampling; 2-approximation for max--min dispersion.

\bibitem{jaccard}
P.~Jaccard.
\newblock The distribution of the flora in the alpine zone.
\newblock \emph{New Phytologist}, 11(2):37--50, 1912.
\newblock \doi{10.1111/j.1469-8137.1912.tb05611.x}

\bibitem{pymatgen}
S.~P.~Ong, W.~D.~Richards, A.~Jain, G.~Hautier, M.~Kocher, S.~Cholia, D.~Gunter, V.~L.~Chevrier,
K.~A.~Persson, and G.~Ceder.
\newblock Python Materials Genomics (pymatgen): A robust, open-source python library for materials
analysis.
\newblock \emph{Computational Materials Science}, 68:314--319, 2013.
\newblock \doi{10.1016/j.commatsci.2012.10.028}
\newblock Used for formula parsing, reduction and conventional element ordering.

\end{thebibliography}

\end{document}
